% ============ 2. Related Literature and the Research Gap ============
\section{Related Literature and the Research Gap}\label{sec:lit}

\subsection{News text and stock returns}

That quantified news text contains information about subsequent returns is one of the longest-standing results in empirical finance. \citet{tetlock2007} measures the pessimism of a daily press column and finds transient market-level pressure and reversal; \citet{garcia2013} shows news sentiment predicts market returns in recessions but not expansions; \citet{chan2003} documents that stocks with salient news drift and then reverse while no-news stocks mean-revert; and \citet{da2011} show that attention itself, measured by search volume, predicts short-run pressure and reversal. The measurement instrument in this first generation is a dictionary~---originally the Harvard psychological lists, then the finance-specific \citet{loughran2011} lists that now dominate the field, and whose construction and properties are reviewed in \citet{loughran2016}. Dictionary measurement is transparent and cheap but blind to negation, context, and event semantics, a limitation already visible in the higher information content of supervised alternatives \citep{jegadeesh2013}.

\subsection{LLMs as measurement instruments}

The current generation replaces the dictionary with an instruction-following LLM. \citet{lopezlira2026} is the reference point: ChatGPT scores of firm-specific headlines predict next-day returns and beat dictionary and FinBERT-style baselines on the same text~---in the published \emph{Journal of Financial Economics} version, 159,137 firm-headline observations on 4,123 U.S.\ stocks over October 2021--May 2024, daily hit rates near 90\% for the initial reaction, a 2.97 Sharpe for the drift strategy, and a Sharpe that decays from 6.54 to 1.22 across the sample as LLM adoption spreads~---the market-efficiency arc our audit's fragility findings rhyme with at small scale. The published version anchors a fast-growing literature surveyed by \citet{nie2024survey}, including LLM sentiment scores as trading signals against dictionary baselines \citep{kirtac2024}, cross-model comparisons \citep{teles2025}, and prompt-engineered GPT-4o versus FinBERT classification \citep{kang2025}. The premise of all of this work is that the LLM's score is a \emph{function of the text}. The present paper interrogates that premise directly.

\subsection{The audit strand: bias, memory, and stability}

A newer strand asks what else the score is a function of. \citet{glasserman2023} are the direct precedent: they compare trading strategies built on original versus \emph{de-identified} headlines inside the model's training window and find, surprisingly, that anonymized headlines outperform~---company identity acts as a ``distraction'' that corrupts sentiment measurement, an effect they show is strongest for large firms. \citet{lopezlira2025} demonstrate that LLMs recall pre-cutoff macroeconomic and financial data with near-lookup accuracy, that instructions and masking fail to suppress recall, and that post-cutoff there is no recall~---which makes post-cutoff windows the clean environment for studying the \emph{remaining} channels. \citet{didisheim2025} show predictable memory in financial analysis, stronger at the index level than the individual-stock level, suggesting firm-level scores are less contaminated by recall. \citet{chen2025} find that LLMs asked to forecast from historical returns extrapolate recent performance~---a mechanism that matters here because web-searched corpora over-represent headlines that \emph{describe} recent performance. On the instrument side, \citet{herrera2025} document the ``model variability problem'': sentiment classifications differ materially across models and across repeated calls of the same model. And standardized look-ahead audits are beginning to appear as benchmarks~---most recently \citet{benhenda2026}'s Look-Ahead-Bench, which scores point-in-time versus standard LLMs by alpha decay across market regimes~---a field-level recognition that output-based evaluation cannot certify measurement validity. Convergence is arriving from the econometrics and evaluation sides as well: \citet{christensen2026} supply formal inference corrections~---with a public package~---for regressions on AI-generated covariates whose measurement error cannot be validated away, which is precisely the situation of a single-scorer sentiment pipeline; \citet{bisharat2026} benchmark twelve sentiment instruments, from dictionaries to open-source LLMs, against evaluation criteria that go beyond return correlation; and the ten-author position paper of \citet{kong2026}~---sharing an author with the \citet{nie2024survey} survey~---names five recurring biases in financial-LLM evaluation (look-ahead, survivorship, narrative, objective, cost), reports that no single one is discussed in more than 28\% of 164 reviewed papers, and demands structural-validity checks before any deployment claim; the audit this paper reports is a within-study instance of that demand, its H4 and H5 the look-ahead and selection entries. Outside finance, the NLP literature has documented context effects that rhyme with the channels this paper quantifies: batching multiple data points into one prompt degrades accuracy and makes outputs sensitive to the position and order of items within the prompt \citep{lin2024batch, liu2024lost}; LLMs carry systematic selection biases over presented options \citep{zheng2024}; and the scores themselves move with rating-scale design~---format, range, sign~---by enough to re-bin a fifth of texts \citep{jo2026}, of which the level channels quantified below are finance-specific instances. What none of these audits measures, and what the NLP results leave unstudied in a measurement setting, is the scoring \emph{context} itself: the batch a headline is scored in, the sitting, and the identity it carries. The present audit adds all three to the checklist.

\subsection{The gap}

Three gaps remain. First, \citet{glasserman2023}'s anonymization experiment is run \emph{in-sample}; whether identity effects survive in a post-cutoff window~---where they can no longer be memory~---has not been tested, and the answer determines how much of the identity effect is a pure prior (about the firm) versus recall (about the outcome). Second, the contamination channels have been studied one at a time; no study measures identity, stochasticity, text-level price echoes, and sample selection \emph{jointly} against the same return windows with formal difference tests and family-wise error control. Third, small-sample LLM predictability studies~---including, in its first pass, this very project~---routinely compare instruments by the significance of one coefficient against another, a practice the econometrics literature rejects \citep[e.g.,][]{romano2005} and that we correct here with bootstrap difference tests plus formal equivalence assessments at stated margins. The contribution of this paper is to fill these three gaps on one frozen corpus, and to convert the resulting checklist into a reusable protocol.

\subsection{Hypotheses}

The audit family, pre-specified before the augmented inference was run:

\begin{description}[leftmargin=1.6em, itemsep=3pt]
\item[H1 (Baseline).] Named LLM sentiment positively predicts the next-day close-to-close return, conditional on lagged return, lagged volatility, and the market factor, under two-way clustered inference, a wild-cluster bootstrap, and Benjamini--Hochberg control across the declared family.
\item[H2 (Identity).] Anonymizing firm identifiers (a) shifts the level of the score distribution and (b) leaves the predictive coefficient statistically unchanged, the two statements together implying an identity \emph{prior} without identity \emph{information}; the named-versus-dictionary difference is likewise tested directly rather than inferred from separate significances. [As originally fielded, part (a) is confounded with session drift; the pre-specified interleaved session of Section~\ref{sec:results:h2} resolves it~---the level shift is $-2.82$ points (round-level $t=6.8$; call-clustered $t=2.7$), re-attributed by the fifth session's swap control to de-identification rather than any firm-specific prior~---and part (b) now carries a formal equivalence verdict at the half-benchmark margin in all four designs (Appendix~\ref{app:power}).]
\item[H3 (Stability).] Repeat calls of the identical prompt are highly but imperfectly correlated; the intra-class correlation quantifies the stochastic component of the instrument and the attenuation it implies.
\item[H4 (Echo).] Headlines containing realized-price language carry the predictive loading; clean headlines do not; and the same-day loading survives anonymization, locating the channel in the text.
\item[H5 (Selection).] Coverage (which firm-days carry news) is positively selected on realized absolute returns and volatility.
\end{description}

H2 and H4 are the hypotheses whose null results would have been hardest to publish a decade ago and are the most informative now: H2 distinguishes priors from information, and H4 distinguishes text semantics from text echoes of prices.
